Supernovae of type Ia are considered very important for the measurements of
large extragalactic distances. The brightening and subsequent dimming of these
explosions follow a characteristic light curve, which helps in identifying
these as supernovae of type Ia.

Light curves of all type Ia supernovae can be fit to the same model light
curve, when they are scaled appropriately. In order to achieve this, we first
have to express the light curves in the reference frame of the host galaxy by
taking care of the cosmological stretching/dilation of all observed time
intervals, $\Delta t_{\mathrm{obs}}$, by a factor of $(1+z)$. The time
interval in the rest frame of the host galaxy is denoted by
$\Delta t_{\mathrm{gal}}$.

The rest frame light curve of a supernova changes by two magnitudes compared
to the peak in a time interval $\Delta t_0$ after the peak. If we further
scale the time intervals by a factor of $s$ (i.e.\
$\Delta t_s = s\,\Delta t_{\mathrm{gal}}$) such that the scaled value of
$\Delta t_0$ is the same for all supernovae, the light curves turn out to have
the same shape. It also turns out that $s$ is related linearly to the absolute
magnitude, $M_{\mathrm{peak}}$, at the peak luminosity for the supernova. That
is, we can write
\[ s = a + b M_{\mathrm{peak}}, \]
where $a$ and $b$ are constants. Knowing the scaling factor, one can determine
absolute magnitudes of supernovae at unknown distances from the above linear
equation.

The table below contains data for three supernovae, including their distance
moduli, $\mu$ (for the first two), their recession speed, $cz$, and their
apparent magnitudes, $m_{\mathrm{obs}}$, at different times. The time
$\Delta t_{\mathrm{obs}} \equiv t - t_{\mathrm{peak}}$ shows number of days
from the date at which the respective supernova reached peak brightness. The
observed magnitudes have already been corrected for interstellar as well as
atmospheric extinction.

\begin{center}
\begin{tabular}{c|c|c|c}
Name & SN2006TD & SN2006IS & SN2005LZ \\
$\mu$ (mag) & 34.27 & 35.64 & \\
$cz$ (km\,s$^{-1}$) & 4515 & 9426 & 12060 \\
\hline
$\Delta t_{\mathrm{obs}}$ (days) & $m_{\mathrm{obs}}$ (mag)
  & $m_{\mathrm{obs}}$ (mag) & $m_{\mathrm{obs}}$ (mag) \\
\hline
$-15.00$ & 19.41 & 18.35 & 20.18 \\
$-10.00$ & 17.48 & 17.26 & 18.79 \\
$-5.00$  & 16.12 & 16.42 & 17.85 \\
0.00     & 15.74 & 16.17 & 17.58 \\
5.00     & 16.06 & 16.41 & 17.72 \\
10.00    & 16.72 & 16.82 & 18.24 \\
15.00    & 17.53 & 17.37 & 18.98 \\
20.00    & 18.08 & 17.91 & 19.62 \\
25.00    & 18.43 & 18.39 & 20.16 \\
30.00    & 18.64 & 18.73 & 20.48 \\
\end{tabular}
\end{center}

\begin{description}
\item[(D3.1)] Compute $\Delta t_{\mathrm{gal}}$ values for all three
  supernovae, and fill them in the given blank boxes in the data tables on the
  BACK side of the Summary Answersheet. On a graph paper, plot the points and
  draw the three light curves in the rest frame (mark your graph as ``D3.1'').
  \hfill[15]
\item[(D3.2)] Take the scaling factor, $s_2$, for the supernova SN2006IS to be
  1.00. Calculate the scaling factors, $s_1$ and $s_3$, for the other two
  supernovae SN2006TD and SN2005LZ, respectively, by calculating $\Delta t_0$
  for them. \hfill[5]
\item[(D3.3)] Compute the scaled time differences, $\Delta t_s$, for all three
  supernovae. Write the values for $\Delta t_s$ in the same data tables on the
  Summary Answersheet. On another graph paper, plot the points and draw the 3
  light curves to verify that they now have an identical profile (mark your
  graph as ``D3.3''). \hfill[14]
\item[(D3.4)] Calculate the absolute magnitudes at peak brightness,
  $M_{\mathrm{peak},1}$, for SN2006TD and $M_{\mathrm{peak},2}$, for SN2006IS.
  Use these values to calculate $a$ and $b$. \hfill[6]
\item[(D3.5)] Calculate the absolute magnitude at peak brightness,
  $M_{\mathrm{peak},3}$, and distance modulus, $\mu_3$, for SN2005LZ.
  \hfill[4]
\item[(D3.6)] Use the distance modulus $\mu_3$ to estimate the value of
  Hubble's constant, $H_0$. Further, estimate the characteristic age of the
  universe, $T_{\mathrm{H}}$. \hfill[6]
\end{description}
